%% ====================================================================
%% Abstract
%% ====================================================================
\begin{abstract}

\textbf{Background.}
The original ARA-Net submission was limited by three major validity problems: weak external classification evidence, an unsupported attention-based Clinical Alignment Score (CAS), and a non-significant Braak-stage analysis. We rebuilt the work as a leakage-aware, subject-level, atlas-guided multimodal Alzheimer's disease (AD) staging study with direct external classification, uncertainty quantification, MCI/AD error analysis, and bounded structural neurodegeneration validation.

\textbf{Methods.}
We constructed subject-level splits across ADNI, AIBL, IXI, and OASIS. ADNI was used for training, validation, and internal testing; AIBL was divided into adaptation training, adaptation validation, and a locked heldout external test; IXI served as a healthy negative-control cohort; and OASIS was retained only as an external stress test. The final algorithm is RC-SPE, a risk-constrained subject-level probability ensemble. RC-SPE combines six base-model probability streams using log-probability pooling, non-negative model weights, class-specific offsets, temperature scaling, and subject-level probability averaging. The previous attention-only CAS claim was removed. Because direct Braak-stage validation was not available, biological evidence was reframed as atlas structural neurodegeneration consistency in a priori AD-relevant regions.

\textbf{Results.}
The original v3 model did not support the earlier external-generalization claim, with AIBL balanced accuracy (BAcc) of 0.399 and IXI healthy CN retention of 0.439. Final subject-level RC-SPE achieved AIBL heldout accuracy of 0.903, BAcc of 0.833, macro AUC of 0.937, AD-vs-CN AUC of 1.000, and CN/MCI/AD recall of 0.961/0.686/0.852. Bootstrap 95\% confidence intervals were 0.759--0.899 for BAcc, 0.894--0.974 for macro AUC, 0.531--0.839 for MCI recall, and 0.710--0.966 for AD recall. IXI healthy CN retention was 1.000. AIBL errors were concentrated near MCI/AD boundaries, with no AD endpoint unit misclassified as CN. OASIS transfer remained weak and is reported as a limitation. The AIBL heldout AD-key atlas-volume consistency score was 0.510 versus a uniform regional null of 0.286, with bootstrap CI 0.479--0.526 and permutation $p=0.026$.

\textbf{Conclusion.}
The revised work supports a domain-adapted, subject-level, atlas-guided multimodal AD staging framework with strong locked AIBL heldout performance, preserved IXI healthy specificity, explicit residual MCI uncertainty, and atlas-region structural neurodegeneration consistency. It does not claim pure zero-shot transfer, solved OASIS generalization, direct Braak-stage validation, or deployment-ready clinical performance.

\end{abstract}

\vspace{0.5em}
\noindent\textbf{Keywords:} Alzheimer's disease; structural MRI; atlas-guided multimodal learning; external validation; subject-level staging; neurodegeneration consistency; open-source research prototype

\vspace{0.5em}
\noindent\textbf{Abbreviations:} AD, Alzheimer's disease; ADNI, Alzheimer's Disease Neuroimaging Initiative; AIBL, Australian Imaging, Biomarkers and Lifestyle study; AUC, area under the receiver-operating-characteristic curve; BAcc, balanced accuracy; CAS, Clinical Alignment Score; CN, cognitively normal; IXI, Information eXtraction from Images dataset; MCI, mild cognitive impairment; MMSE, Mini-Mental State Examination; OASIS, Open Access Series of Imaging Studies; sMRI, structural magnetic resonance imaging.


%% ====================================================================
\section{Introduction}
\label{sec:intro}
%% ====================================================================

Structural MRI is widely used in AD research because it captures neurodegeneration patterns including medial temporal atrophy, ventricular enlargement, and broader brain-volume change~\citep{jack2018nia,frisoni2010neuroimaging}. These anatomical changes are not sufficient for clinical diagnosis by themselves, but they provide a biologically meaningful substrate for machine-learning models that aim to support CN/MCI/AD staging. A credible staging model must therefore demonstrate not only within-cohort accuracy, but also leakage-aware external evaluation, class-specific error behavior, and biological interpretation that is proportional to the evidence actually available.

The original ARA-Net manuscript attempted to connect atlas-guided attention with external generalization and neuropathological plausibility. After reviewer evaluation, that framing was no longer defensible. First, cross-dataset generalization was inferred primarily from explanation similarity rather than from direct external CN/MCI/AD classification. Second, the proposed Clinical Alignment Score did not provide valid support for the claimed alignment, because the attention-derived result was weaker than the intended biological interpretation. Third, the Braak-stage analysis was non-significant, so direct pathological-stage claims could not be retained.

The revised work therefore changes the scientific question. Instead of asking whether attention maps alone can validate ARA-Net, we ask whether an atlas-guided multimodal framework can support subject-level external AD staging after explicit external-domain adaptation, while preserving healthy negative-control specificity and reporting remaining weaknesses. The final model target is not an attention score. It is locked AIBL heldout subject-level CN/MCI/AD classification with bootstrap uncertainty and transparent MCI/AD error analysis.

This change also narrows the biological claim. Because direct Braak-stage validation is unavailable, we no longer claim pathological staging. Instead, we test whether atlas structural MRI changes concentrate in a priori AD-relevant regions. This analysis is a structural neurodegeneration consistency check, not a proof that model attention identifies biomarkers. The claim is deliberately bounded because the available data support an MRI proxy but not direct neuropathological validation.

The revised contributions are:
\begin{enumerate}[leftmargin=*,itemsep=2pt,label=\textbf{C\arabic*.}]
    \item A leakage-aware subject-level protocol across ADNI, AIBL, IXI, and OASIS.
    \item Direct external CN/MCI/AD classification on a locked AIBL heldout endpoint.
    \item IXI healthy negative-control evaluation to quantify false impairment predictions.
    \item RC-SPE, a risk-constrained subject-level probability ensemble with non-negative model weights, class offsets, temperature scaling, and subject-level probability averaging.
    \item Bootstrap stability analysis and subject-level MCI/AD error analysis.
    \item Replacement of the unsupported attention-only CAS/Braak claim with atlas structural neurodegeneration consistency analysis.
    \item An open-source research deployment package with explicit clinical-use boundaries.
\end{enumerate}

Figure~\ref{fig:v6_framework} summarizes the revised work. The visual story should be unmistakable: this is no longer a cosmetic repair to an attention manuscript. It is a substantially rebuilt external subject-level AD staging study with a bounded biological validation target.

\begin{figure}[p]
    \centering
    \fbox{\begin{minipage}[c][0.63\textheight][c]{0.94\textwidth}
    \textbf{Placeholder for Figure 1. Revised study framework and central model objective.}

    \vspace{0.8em}
    Panels to draw: \textbf{(a)} old v3 failure points and V6 rebuild target; \textbf{(b)} subject-level cohort construction; \textbf{(c)} atlas-guided multimodal features; \textbf{(d)} RC-SPE risk-constrained probability ensemble; \textbf{(e)} endpoint and claim-boundary summary. The visual emphasis should be that the paper is now an external subject-level AD staging framework, not an attention-only interpretability patch.
    \end{minipage}}
    \caption{\textbf{Revised ARA-Net V6 framework.} The revised workflow uses leakage-aware subject-level cohort construction, atlas-guided multimodal feature extraction, RC-SPE risk-constrained probability pooling, and subject-level probability averaging. The primary endpoint is locked AIBL heldout CN/MCI/AD staging; IXI is a healthy negative-control specificity endpoint; OASIS is retained as a stress-test limitation. The original attention-only CAS and direct Braak-stage claims are replaced by atlas structural neurodegeneration consistency analysis and explicit claim-boundary language.}
    \label{fig:v6_framework}
\end{figure}


%% ====================================================================
\section{Methods}
\label{sec:methods}
%% ====================================================================

\subsection{Study design and notation}
\label{sec:study_design}

The revised study uses a multi-cohort, subject-level design. Each scan-level record is represented as
\begin{equation}
r_i = \left(d_i, s_i, t_i, y_i, x_i\right),
\label{eq:record}
\end{equation}
where $d_i$ is the dataset, $s_i$ is the subject identifier, $t_i$ is the visit or scan identifier, $y_i \in \{\mathrm{CN},\mathrm{MCI},\mathrm{AD}\}$ is the diagnostic label when available, and $x_i$ is the multimodal feature vector. The full data inventory is
\begin{equation}
\mathcal{D} =
\mathcal{D}_{\mathrm{ADNI}} \cup
\mathcal{D}_{\mathrm{AIBL}} \cup
\mathcal{D}_{\mathrm{IXI}} \cup
\mathcal{D}_{\mathrm{OASIS}} .
\label{eq:data_union}
\end{equation}

The development and evaluation roles were fixed before final endpoint reporting:
\begin{equation}
\mathcal{D}_{\mathrm{ADNI}} =
\mathcal{D}_{\mathrm{train}} \cup
\mathcal{D}_{\mathrm{val}} \cup
\mathcal{D}_{\mathrm{internal}},
\qquad
\mathcal{D}_{\mathrm{AIBL}} =
\mathcal{D}_{\mathrm{adapt\mbox{-}train}} \cup
\mathcal{D}_{\mathrm{adapt\mbox{-}val}} \cup
\mathcal{D}_{\mathrm{heldout}} .
\label{eq:split_roles}
\end{equation}
IXI was used as a healthy negative-control cohort, and OASIS was held out as an external stress test. OASIS was not used for final model tuning.

Subject-level leakage control was imposed by requiring disjoint subject identifiers between development and evaluation partitions:
\begin{equation}
\mathcal{S}_a \cap \mathcal{S}_b = \varnothing
\quad
\forall a \neq b,\quad
a,b \in
\{\mathrm{train},\mathrm{val},\mathrm{internal},\mathrm{adapt\mbox{-}train},
\mathrm{adapt\mbox{-}val},\mathrm{heldout}\},
\label{eq:subject_disjoint}
\end{equation}
where $\mathcal{S}_a=\{s_i:r_i\in\mathcal{D}_a\}$. This rule is important because repeated imaging from the same participant can otherwise inflate performance estimates~\citep{varoquaux2017cross}.

\subsection{Cohort construction and endpoint inventory}
\label{sec:data}

Table~\ref{tab:v6_splits} summarizes the split inventory. The subject counts in the table are unique split-level participants used to define leakage-free splits. For longitudinal cohorts, endpoint units can differ from unique participants because a participant may contribute repeated scans or distinct diagnostic-state endpoint units over time. The primary endpoint is therefore described explicitly as subject-level diagnostic-state classification after probability aggregation.

\begin{table}[t]
\centering
\caption{V6 split inventory. Counts refer to scan inventory and unique split-level subjects.}
\label{tab:v6_splits}
\small
\begin{tabular}{@{}lrrr@{}}
\toprule
Split & Scans & Subjects & Role \\
\midrule
ADNI train & 1{,}686 & 450 & Main training \\
ADNI validation & 355 & 97 & Validation and calibration \\
ADNI internal test & 360 & 96 & Internal risk analysis \\
AIBL adaptation training & 719 & 385 & External-domain adaptation \\
AIBL adaptation validation & 191 & 105 & External-domain calibration \\
AIBL locked heldout & 397 & 210 & Primary external test \\
OASIS external & 99 & 99 & Stress test \\
IXI healthy external & 581 & 581 & Healthy negative-control specificity \\
\bottomrule
\end{tabular}
\end{table}

For a subject-level diagnostic-state unit $u$, scan indices were grouped as
\begin{equation}
\mathcal{I}_u =
\{i: s_i=s(u),\; y_i=y(u),\; d_i=d(u),\; r_i\in\mathcal{D}_{q(u)}\},
\label{eq:endpoint_group}
\end{equation}
where $q(u)$ is the split assignment. This grouping prevents scan-level repetitions from being interpreted as independent final endpoint units. Scan-level results are retained only as reference analyses.

\begin{figure}[p]
    \centering
    \fbox{\begin{minipage}[c][0.60\textheight][c]{0.94\textwidth}
    \textbf{Placeholder for Figure 2. Data line and endpoint design.}

    \vspace{0.8em}
    Panels to draw: \textbf{(a)} ADNI/AIBL/OASIS/IXI cohort blocks with scan and subject counts; \textbf{(b)} ADNI train--validation--internal-test split; \textbf{(c)} AIBL adaptation train/validation and locked heldout split; \textbf{(d)} endpoint aggregation from scans to subject-level diagnostic-state units; \textbf{(e)} OASIS explicitly shown as stress test, not tuning or success validation.
    \end{minipage}}
    \caption{\textbf{Leakage-aware multi-cohort data line.} The revised study separates model development, external-domain adaptation, locked external testing, healthy-control specificity testing, and stress testing. This figure should make clear that AIBL heldout is the primary external endpoint, IXI evaluates false impairment in healthy controls, and OASIS is preserved only as a limitation stress test.}
    \label{fig:v6_data_line}
\end{figure}

\subsection{Atlas-guided multimodal feature representation}
\label{sec:features}

MRI features were extracted from a 21-region atlas and included regional volumetric and intensity summaries. The atlas representation is denoted
\begin{equation}
a_i = \left[a_{i,1},a_{i,2},\ldots,a_{i,p_a}\right] \in \mathbb{R}^{p_a}.
\label{eq:atlas_features}
\end{equation}
The core clinical feature vector is denoted
\begin{equation}
c_i = \left[\mathrm{age}_i,\mathrm{sex}_i,\mathrm{education}_i,\mathrm{APOE4}_i,\mathrm{MMSE}_i,\mathrm{CDR\mbox{-}SB}_i\right] \in \mathbb{R}^{p_c},
\label{eq:clinical_features}
\end{equation}
where unavailable variables were imputed within the training/adaptation workflow and accompanied by missingness indicators when required. The final multimodal representation was
\begin{equation}
x_i = \left[\tilde{a}_i,\tilde{c}_i,m_i\right] \in \mathbb{R}^{p},
\label{eq:feature_concat}
\end{equation}
where $\tilde{a}_i$ and $\tilde{c}_i$ are standardized atlas and clinical features, and $m_i$ contains missingness indicators.

Feature standardization used development-set statistics only:
\begin{equation}
\tilde{x}_{i,j} =
\frac{x_{i,j}-\mu_j^{(\mathrm{dev})}}{\sigma_j^{(\mathrm{dev})}+\eta},
\label{eq:standardization}
\end{equation}
where $\eta$ is a small numerical constant. The development statistics were not recomputed on locked heldout endpoint labels.

The revised model should therefore be described as atlas-guided and multimodal. The atlas component anchors the model to interpretable neuroanatomical regions, while the clinical component reflects the known diagnostic signal contained in age, global cognition, and functional staging. The central evidence is classification performance and structural consistency, not attention-map biomarker discovery.

\subsection{Candidate models and comparator roles}
\label{sec:candidates}

The revised experimental framework included atlas-only, cascade, atlas+clinical, clinical-only, biomarker-enhanced, and ensemble models. These models have different interpretive roles. Atlas-only and cascade models test whether MRI-derived regional structure alone can provide sufficient external staging. Atlas+clinical models test the main scientific target: anatomically grounded multimodal staging. Clinical-only models act as upper-bound comparators because clinical variables can encode strong diagnostic signal but do not retain the MRI atlas-guided mechanism. Biomarker-enhanced variants provide sensitivity analyses when additional variables are available.

For each base model $m$, class probabilities were represented as
\begin{equation}
p_m(y=k\mid x_i)=
\frac{\exp(f_{m,k}(x_i))}
{\sum_{c\in\mathcal{Y}}\exp(f_{m,c}(x_i))},
\qquad
\mathcal{Y}=\{\mathrm{CN},\mathrm{MCI},\mathrm{AD}\}.
\label{eq:base_prob}
\end{equation}
When a model family produced calibrated probabilities directly, $f_{m,k}$ denotes its logit-scale representation. When a tree-based model produced probability estimates, the probabilities were clipped before log-pooling to avoid numerical underflow.

\subsection{RC-SPE: risk-constrained subject-level probability ensemble}
\label{sec:model}

The final locked algorithm was RC-SPE, a risk-constrained subject-level probability ensemble. RC-SPE was designed to avoid a common failure mode in external AD staging: improving aggregate accuracy while sacrificing minority disease recall or overcalling impairment in healthy external controls. It combined six base-model probability streams:
\begin{itemize}[leftmargin=*,itemsep=2pt]
    \item AIBL-adapted atlas-biomarker-enhanced HGB.
    \item AIBL-adapted atlas-core-clinical HGB.
    \item AIBL-adapted clinical-biomarker-only RF balanced.
    \item AIBL-adapted clinical-core-only HGB.
    \item AIBL-adapted clinical-core-only RF balanced.
    \item RF--logistic regression ensemble component.
\end{itemize}

Let $p_{m,k}(x_i)$ denote the class probability for scan $i$, class $k$, and base model $m$. The final ensemble uses log-probability pooling:
\begin{equation}
z_{i,k} =
\frac{1}{T}\sum_{m=1}^{M} w_m
\log\left(\max(p_{m,k}(x_i), \epsilon)\right) + b_k ,
\qquad
k \in \{\mathrm{CN},\mathrm{MCI},\mathrm{AD}\},
\label{eq:v6_log_pool}
\end{equation}
where
\begin{equation}
w_m \geq 0,\qquad \sum_{m=1}^{M}w_m=1,\qquad T>0.
\label{eq:ensemble_constraints}
\end{equation}
Here $w_m$ are non-negative model weights, $T$ is a temperature parameter, $b_k$ is a class-specific offset, and $\epsilon$ prevents taking the logarithm of zero. Calibrated class probabilities are:
\begin{equation}
\tilde{p}_{i,k} =
\frac{\exp(z_{i,k})}
{\sum_{c\in\mathcal{Y}}\exp(z_{i,c})}.
\label{eq:v6_softmax}
\end{equation}

RC-SPE was selected using ADNI validation, AIBL adaptation validation, and IXI healthy-control specificity, with OASIS excluded from the final tuning objective. In compact form, the selection objective can be written as
\begin{align}
\max_{\theta}
\quad
\mathcal{J}(\theta)
&=
\alpha_1 \operatorname{BAcc}_{\mathrm{ADNI\ val}}(\theta)
+ \alpha_2 \operatorname{AUC}_{\mathrm{ADNI\ val}}(\theta)
+ \alpha_3 \operatorname{BAcc}_{\mathrm{AIBL\ adapt\ val}}(\theta) \notag\\
&\quad
+ \alpha_4
\min\{
\operatorname{Rec}^{\mathrm{AIBL\ adapt\ val}}_{\mathrm{MCI}}(\theta),
\operatorname{Rec}^{\mathrm{AIBL\ adapt\ val}}_{\mathrm{AD}}(\theta)
\}
+ \alpha_5
\operatorname{Ret}^{\mathrm{IXI}}_{\mathrm{CN}}(\theta),
\label{eq:selection_objective}
\end{align}
where $\theta=(w,b,T)$ and the $\alpha$ terms define the balanced rescue profile. This objective reflects the practical aim of rescuing minority-class performance while preserving healthy-control specificity.

Because IXI contains healthy subjects only, we also tracked false impairment as a hard risk indicator:
\begin{equation}
\operatorname{Risk}_{\mathrm{IXI}}(\theta)
=
1-\operatorname{Ret}^{\mathrm{IXI}}_{\mathrm{CN}}(\theta),
\label{eq:ixi_risk}
\end{equation}
and treated candidate profiles with improved AIBL MCI recall but increased IXI false impairment as riskier alternatives rather than as the locked final model.

\begin{center}
\fbox{\begin{minipage}{0.92\textwidth}
\textbf{Algorithm 1. RC-SPE for locked subject-level AD staging.}
\begin{enumerate}[leftmargin=*,itemsep=1pt]
    \item Train or load $M$ base probability models and obtain $p_{m,k}(x_i)$ for each scan.
    \item Select non-negative weights $w_m$, offsets $b_k$, and temperature $T$ on ADNI validation, AIBL adaptation validation, and IXI healthy specificity; do not use OASIS for tuning.
    \item Pool probabilities by Eq.~\ref{eq:v6_log_pool} and transform logits by Eq.~\ref{eq:v6_softmax}.
    \item Aggregate repeated scans into subject-level endpoint probabilities by Eq.~\ref{eq:v6_subject_average}.
    \item Lock $\theta=(w,b,T)$ and report AIBL heldout, IXI healthy specificity, OASIS stress-test, bootstrap, calibration, and error analyses.
\end{enumerate}
\end{minipage}}
\end{center}

\begin{figure}[p]
    \centering
    \fbox{\begin{minipage}[c][0.60\textheight][c]{0.94\textwidth}
    \textbf{Placeholder for Figure 3. RC-SPE locked subject-level probability ensemble.}

    \vspace{0.8em}
    Panels to draw: \textbf{(a)} six base probability streams; \textbf{(b)} RC-SPE objective with AIBL minority recall and IXI healthy specificity; \textbf{(c)} log-probability pooling with weights, class offsets, and temperature; \textbf{(d)} scan-level to subject-level probability averaging; \textbf{(e)} final CN/MCI/AD output with confidence and margin. The main visual message is that the method is a risk-constrained calibrated subject-level AD staging algorithm, not a generic average ensemble.
    \end{minipage}}
    \caption{\textbf{RC-SPE locked subject-level probability ensemble.} The final algorithm combines six base-model probability streams through log-probability pooling, class-specific offsets, temperature scaling, and subject-level probability averaging. Its selection objective balances external AIBL staging performance, minority-class rescue, and IXI healthy-control specificity while excluding OASIS from tuning.}
    \label{fig:v6_model}
\end{figure}

\subsection{Subject-level probability aggregation}
\label{sec:subject_aggregation}

For each subject-level endpoint unit $u$, repeated scan probabilities were averaged:
\begin{equation}
\bar{p}_{u,k} =
\frac{1}{|\mathcal{I}_u|}
\sum_{i\in\mathcal{I}_u}\tilde{p}_{i,k},
\qquad
\hat{y}_u = \arg\max_{k\in\mathcal{Y}}\bar{p}_{u,k}.
\label{eq:v6_subject_average}
\end{equation}
Prediction confidence and decision margin were defined as
\begin{equation}
\operatorname{conf}(u)=\max_k \bar{p}_{u,k},
\qquad
\operatorname{margin}(u)=
\max_k \bar{p}_{u,k}
-
\max_{c\neq \hat{y}_u}\bar{p}_{u,c}.
\label{eq:confidence_margin}
\end{equation}
These quantities were used for error analysis but not to change locked predictions.

\subsection{Classification metrics}
\label{sec:metrics}

Let $C_{ab}$ be the confusion-matrix entry for true class $a$ and predicted class $b$. Per-class recall, precision, and balanced accuracy were defined as
\begin{equation}
\operatorname{Rec}_a =
\frac{C_{aa}}{\sum_b C_{ab}},
\qquad
\operatorname{Prec}_a =
\frac{C_{aa}}{\sum_b C_{ba}},
\qquad
\operatorname{BAcc} =
\frac{1}{|\mathcal{Y}_{\mathrm{present}}|}
\sum_{a\in\mathcal{Y}_{\mathrm{present}}}\operatorname{Rec}_a .
\label{eq:metrics}
\end{equation}
Overall accuracy was
\begin{equation}
\operatorname{Acc}=
\frac{\sum_a C_{aa}}{\sum_a\sum_b C_{ab}}.
\label{eq:accuracy}
\end{equation}
For multiclass AUC, one-vs-rest AUCs were averaged across evaluable classes:
\begin{equation}
\operatorname{MacroAUC}=
\frac{1}{|\mathcal{Y}_{\mathrm{AUC}}|}
\sum_{k\in\mathcal{Y}_{\mathrm{AUC}}}
\operatorname{AUC}
\left(\mathbf{1}[y=k],\bar{p}_{k}\right).
\label{eq:macro_auc}
\end{equation}
For IXI, all subjects are healthy controls, so the primary specificity metric was CN retention:
\begin{equation}
\operatorname{Ret}^{\mathrm{IXI}}_{\mathrm{CN}}
=
\frac{\#\{u\in\mathrm{IXI}:\hat{y}_u=\mathrm{CN}\}}
{\#\{u\in\mathrm{IXI}\}},
\qquad
\operatorname{FIR}^{\mathrm{IXI}}=1-\operatorname{Ret}^{\mathrm{IXI}}_{\mathrm{CN}},
\label{eq:ixi_retention}
\end{equation}
where FIR denotes the false impairment rate.

\subsection{Bootstrap uncertainty}
\label{sec:bootstrap}

Uncertainty for the locked AIBL heldout subject-level endpoint was estimated using 2{,}000 bootstrap resamples. For a metric $M$ and bootstrap sample $b$, the empirical estimate was
\begin{equation}
\hat{M}^{(b)}
=
M\left(\{(y_u,\bar{p}_u):u\in\mathcal{B}^{(b)}\}\right),
\qquad
\mathcal{B}^{(b)}
\sim
\operatorname{SampleWithReplacement}(\mathcal{U}_{\mathrm{AIBL\ heldout}}).
\label{eq:bootstrap}
\end{equation}
The percentile confidence interval was
\begin{equation}
\widehat{\mathrm{CI}}_{95\%}(M)
=
\left[
Q_{0.025}\{\hat{M}^{(b)}\}_{b=1}^{B},
Q_{0.975}\{\hat{M}^{(b)}\}_{b=1}^{B}
\right],
\qquad B=2{,}000 .
\label{eq:bootstrap_ci}
\end{equation}

\subsection{MCI/AD error analysis}
\label{sec:error_methods}

Subject-level errors were summarized using transition rates:
\begin{equation}
R_{ab}=
\frac{C_{ab}}{\sum_c C_{ac}},
\qquad
a,b\in\{\mathrm{CN},\mathrm{MCI},\mathrm{AD}\}.
\label{eq:transition_rates}
\end{equation}
We also defined error groups for adjacent and severe mistakes:
\begin{equation}
\mathcal{E}_{a\rightarrow b}
=
\{u:y_u=a,\hat{y}_u=b,a\neq b\}.
\label{eq:error_groups}
\end{equation}
Group profiles were summarized for age, MMSE, CDR-SB, APOE4, atlas hippocampal volume, atlas lateral ventricular volume, AD-like atlas z-score, maximum probability, and decision margin. The purpose was to distinguish boundary-like uncertainty from clinically severe collapse of disease cases into CN.

\subsection{Structural neurodegeneration consistency}
\label{sec:bio}

The old attention-centered CAS was removed because it did not provide valid biological evidence. The revised biological analysis tests whether disease-associated atlas volume differences concentrate in a priori AD-relevant regions: bilateral hippocampus, bilateral amygdala, and bilateral lateral ventricles. Let $\mathcal{R}$ be the atlas-region set and $\mathcal{K}\subset\mathcal{R}$ be the six AD-key regions. For each region $r$, the absolute disease-associated volume contrast was
\begin{equation}
\Delta_r =
\left|
\mu_{r,\mathrm{AD}}-\mu_{r,\mathrm{CN}}
\right|,
\label{eq:region_delta}
\end{equation}
where the means were computed within the relevant evaluation group. The AD-key structural consistency score was
\begin{equation}
\operatorname{NCS}_{\mathcal{K}}
=
\frac{\sum_{r\in\mathcal{K}}\Delta_r}
{\sum_{r\in\mathcal{R}}\Delta_r}.
\label{eq:ncs}
\end{equation}
The uniform regional null was
\begin{equation}
\operatorname{NCS}_{0}=
\frac{|\mathcal{K}|}{|\mathcal{R}|}
=
\frac{6}{21}=0.286.
\label{eq:ncs_null}
\end{equation}

Permutation testing evaluated whether the observed concentration exceeded random region-set selection:
\begin{equation}
p_{\mathrm{perm}} =
\frac{1+\sum_{b=1}^{B}
\mathbf{1}
\left[
\operatorname{NCS}_{\mathcal{K}^{(b)}}
\geq
\operatorname{NCS}_{\mathcal{K}}
\right]}
{B+1},
\qquad
|\mathcal{K}^{(b)}|=|\mathcal{K}|.
\label{eq:permutation}
\end{equation}
This is a structural MRI neurodegeneration proxy. It does not provide postmortem-stage validation or attention-map biomarker evidence.

\subsection{Open-source research deployment and clinical-use boundary}
\label{sec:deployment}

The public implementation includes reproducible analysis scripts, a research inference wrapper, an HTTP API/static frontend, aggregate reports, figures, documentation, and toy probability examples. Raw ADNI, AIBL, OASIS, and IXI data are governed by their original access agreements and are not redistributed. Row-level private predictions, clinical spreadsheets, MRI volumes, and non-redistributable checkpoints are excluded from the public release.

The software is released as an open-source research prototype for retrospective evaluation and future prospective validation. It is not a medical device, is not cleared or approved for clinical use, and is not intended for standalone diagnosis or patient-care decisions.


%% ====================================================================
\section{Results}
\label{sec:results}
%% ====================================================================

\subsection{External classification failure in the original v3 model}
\label{sec:old_failure}

The original v3 model did not support the earlier cross-dataset generalization claim. On AIBL, it achieved accuracy 0.606, BAcc 0.399, and macro AUC 0.597. On IXI, only 0.439 of healthy controls were retained as CN, indicating a high false impairment rate. These results motivate the central revision: external validation must be based on direct classification outcomes and healthy-control specificity, not on explanation similarity alone.

This failure also explains why the final manuscript cannot be framed as a small textual correction. The model, endpoint, claim boundary, and biological analysis were all changed. The revised evidence focuses on locked AIBL subject-level performance, IXI specificity, bootstrap stability, and bounded structural consistency.

\subsection{Locked AIBL heldout subject-level performance}
\label{sec:external_results}

Final locked RC-SPE achieved strong AIBL heldout subject-level performance. It achieved accuracy 0.903, BAcc 0.833, macro AUC 0.937, AD-vs-CN AUC 1.000, and CN/MCI/AD recall 0.961/0.686/0.852. The scan-level reference result was similar: accuracy 0.909, BAcc 0.820, macro AUC 0.939, AD-vs-CN AUC 0.998, and CN/MCI/AD recall 0.964/0.642/0.854.

IXI healthy negative-control specificity was preserved. The final model retained all IXI healthy subjects as CN at both scan and subject levels, corresponding to CN retention 1.000 and false impairment rate 0.000. This result does not provide AD/MCI staging evidence because IXI contains healthy controls only, but it directly addresses whether the final model overcalls impairment in a healthy external cohort.

\begin{table}[t]
\centering
\caption{Main external classification results. The locked primary endpoint is final RC-SPE at subject level on AIBL heldout.}
\label{tab:v6_classification}
\scriptsize
\resizebox{\textwidth}{!}{%
\begin{tabular}{@{}lllrrrrlll@{}}
\toprule
Model/protocol & Unit & Test cohort & Endpoint $n$ & Acc & BAcc & Macro AUC & AD-vs-CN AUC/CN retention & CN/MCI/AD recall & Role \\
\midrule
Old v3 ensemble & scan & AIBL external & 1{,}307 & 0.606 & 0.399 & 0.597 & NA & NA & Failed external baseline \\
Old v3 ensemble & scan & IXI healthy & 581 & 0.439 & 0.439 & NA & CN retention 0.439 & 0.439/0.000/0.000 & Failed specificity baseline \\
v4 atlas+clinical HGB & scan & AIBL heldout & 397 & 0.882 & 0.741 & 0.942 & AD-vs-CN AUC 0.990 & 0.964/0.528/0.732 & Earlier rebuild \\
Final RC-SPE & scan & AIBL heldout & 397 & 0.909 & 0.820 & 0.939 & AD-vs-CN AUC 0.998 & 0.964/0.642/0.854 & Scan-level reference \\
Final RC-SPE & subject & AIBL heldout & 216 & 0.903 & 0.833 & 0.937 & AD-vs-CN AUC 1.000 & 0.961/0.686/0.852 & Locked primary result \\
Final RC-SPE & subject & IXI healthy & 581 & 1.000 & 1.000 & NA & CN retention 1.000 & 1.000/0.000/0.000 & Specificity check \\
Final RC-SPE & subject & OASIS stress & 99 & 0.586 & 0.334 & 0.554 & AD-vs-CN AUC 0.371 & 0.966/0.034/0.000 & Limitation \\
Clinical-only RF comparator & scan & AIBL heldout & 397 & 0.922 & 0.835 & 0.957 & AD-vs-CN AUC 0.997 & 0.970/0.755/0.780 & Upper-bound comparator \\
\bottomrule
\end{tabular}%
}
\end{table}

\begin{figure}[p]
    \centering
    \fbox{\begin{minipage}[c][0.62\textheight][c]{0.94\textwidth}
    \textbf{Placeholder for Figure 4. External classification performance.}

    \vspace{0.8em}
    Panels to draw: \textbf{(a)} AIBL BAcc comparison: old v3, v4 atlas+clinical, final scan, final subject, clinical-only comparator; \textbf{(b)} AIBL macro AUC comparison; \textbf{(c)} final subject-level CN/MCI/AD recall; \textbf{(d)} IXI CN retention and false impairment rate; \textbf{(e)} OASIS as a small stress-test limitation panel. The central story is old external failure to new locked AIBL subject-level success with IXI specificity.
    \end{minipage}}
    \caption{\textbf{External classification performance.} The original v3 model failed external classification and healthy specificity. Final subject-level RC-SPE substantially improves locked AIBL heldout CN/MCI/AD staging and preserves IXI healthy specificity. OASIS is shown as a stress-test limitation, not as a successful validation cohort.}
    \label{fig:v6_external_performance}
\end{figure}

\subsection{Comparator interpretation}
\label{sec:comparators}

The clinical-only RF comparator achieved AIBL heldout BAcc 0.835 and CN/MCI/AD recall 0.970/0.755/0.780 at scan level. This model is reported as a comparator and upper bound rather than the central ARA-Net model, because it does not retain the atlas-guided MRI component central to the revised scientific objective. This comparison is important: it shows that clinical variables carry substantial diagnostic signal and prevents overstating MRI-only interpretability.

RC-SPE is therefore not presented as the single numerically best classifier under every possible feature setting. It is presented as the final model that balances the paper's core requirements: atlas-guided multimodal staging, high locked AIBL subject-level performance, MCI/AD rescue, and IXI healthy specificity. This distinction should remain explicit in the manuscript and figures.

\subsection{Algorithmic ablation and calibration of RC-SPE}
\label{sec:algorithm_results}

We evaluated RC-SPE against single-model, simple averaging, equal log-pooling, partial-parameter, and aggregation ablations using the same private prediction streams. Only aggregate metrics are reported. The ablation demonstrates that the final locked method is not equivalent to a generic average ensemble. The best single base model achieved AIBL BAcc 0.756, MCI recall 0.571, AD recall 0.741, two AD-to-CN errors, and IXI CN retention 0.997. Equal log-pooling without learned weights, offsets, or temperature achieved AIBL BAcc 0.648 and MCI recall 0.171. By contrast, full RC-SPE at subject level achieved AIBL BAcc 0.833, MCI recall 0.686, AD recall 0.852, zero AD-to-CN errors, and IXI CN retention 1.000.

\begin{table}[t]
\centering
\caption{Algorithmic ablation of RC-SPE. OASIS is reported only as a stress-test metric and was not used for final model selection.}
\label{tab:rcspe_ablation}
\scriptsize
\resizebox{\textwidth}{!}{%
\begin{tabular}{@{}lrrrrrrrr@{}}
\toprule
Variant & AIBL BAcc & AIBL AUC & MCI recall & AD recall & AD$\rightarrow$CN & IXI CN retention & ECE & OASIS BAcc \\
\midrule
Best single base model & 0.756 & 0.945 & 0.571 & 0.741 & 2 & 0.997 & 0.122 & 0.310 \\
Arithmetic mean ensemble & 0.711 & 0.951 & 0.400 & 0.741 & 3 & 1.000 & 0.112 & 0.333 \\
Equal log-pooling & 0.648 & 0.954 & 0.171 & 0.778 & 5 & 1.000 & 0.094 & 0.333 \\
Final weights only & 0.815 & 0.955 & 0.657 & 0.815 & 0 & 1.000 & 0.136 & 0.333 \\
Weights + offsets & 0.823 & 0.926 & 0.657 & 0.852 & 0 & 1.000 & 0.219 & 0.310 \\
Weights + temperature & 0.815 & 0.954 & 0.657 & 0.815 & 0 & 1.000 & 0.052 & 0.333 \\
Full RC-SPE, scan-level & 0.810 & 0.940 & 0.660 & 0.805 & 0 & 1.000 & 0.081 & 0.334 \\
Full RC-SPE, subject-level & 0.833 & 0.937 & 0.686 & 0.852 & 0 & 1.000 & 0.078 & 0.334 \\
MCI-rescue profile & 0.811 & 0.952 & 0.886 & 0.593 & 0 & 0.959 & 0.156 & 0.367 \\
AD-rescue profile & 0.698 & 0.930 & 0.143 & 0.963 & 0 & 1.000 & 0.293 & 0.328 \\
\bottomrule
\end{tabular}%
}
\end{table}

The calibration analysis further supports the temperature-scaled component of RC-SPE. The best single base model had AIBL expected calibration error (ECE) 0.122 and negative log-likelihood (NLL) 0.416. The final RC-SPE reduced ECE to 0.078 and NLL to 0.320 while improving BAcc. The weights+temperature ablation achieved the lowest ECE (0.052) and NLL (0.298), but its AIBL BAcc remained 0.815 and AD recall 0.815. The locked final RC-SPE therefore prioritizes the combined endpoint of external staging, AD/MCI rescue, zero AD-to-CN errors, and IXI specificity rather than calibration alone.

We also evaluated risk-profile alternatives. A high-MCI-recall profile increased AIBL MCI recall to 0.886 but reduced AD recall to 0.593 and IXI CN retention to 0.959, meaning 24 IXI healthy subjects were predicted as impaired. A high-AD-recall profile increased AD recall to 0.963 but reduced MCI recall to 0.143. These tradeoffs support the risk-constrained final selection. Leave-one-model-out sensitivity remained stable: dropping any one of the six base streams produced AIBL BAcc 0.823--0.835 and preserved zero AD-to-CN errors, indicating that RC-SPE is not dependent on a single fragile base model.

\subsection{Subject-level MCI and AD error analysis}
\label{sec:error_results}

On the locked AIBL heldout subject-level set, errors were concentrated at disease-stage boundaries. Among 154 CN endpoint units, 148 were classified as CN, five as MCI, and one as AD. Among 35 MCI endpoint units, 24 were classified as MCI, two as CN, and nine as AD. Among 27 AD endpoint units, 23 were classified as AD and four as MCI; no AD endpoint unit was misclassified as CN.

\begin{table}[t]
\centering
\caption{AIBL heldout subject-level confusion matrix for final RC-SPE.}
\label{tab:v6_confusion}
\begin{tabular}{@{}lrrrr@{}}
\toprule
True label & Predicted CN & Predicted MCI & Predicted AD & Recall \\
\midrule
CN & 148 & 5 & 1 & 0.961 \\
MCI & 2 & 24 & 9 & 0.686 \\
AD & 0 & 4 & 23 & 0.852 \\
\bottomrule
\end{tabular}
\end{table}

The transition rates reinforce the same interpretation. MCI-to-AD errors accounted for 0.257 of MCI endpoint units, whereas MCI-to-CN errors accounted for 0.057. AD-to-MCI errors accounted for 0.148 of AD endpoint units, and AD-to-CN errors were 0.000. This pattern is more clinically plausible than a classifier that preserves global accuracy by assigning most uncertain disease cases to CN.

The feature-profile analysis supported a boundary-error interpretation. AIBL AD endpoint units correctly classified as AD had lower MMSE, larger lateral ventricular volume, and higher AD-like atlas z-scores than AD endpoint units classified as CN/MCI. MCI endpoint units classified as AD had lower MMSE and more AD-like atlas profiles than MCI endpoint units classified correctly, consistent with disease-severity boundary uncertainty rather than arbitrary failure.

Internal subject-level performance remained weaker. Internal subject-level BAcc was 0.448 with CN/MCI/AD recall 0.241/0.553/0.550. This internal calibration tension should not be hidden; it is reported as a limitation and motivates future site-specific calibration work. The key external claim remains the locked AIBL heldout endpoint with IXI specificity.

\begin{figure}[p]
    \centering
    \fbox{\begin{minipage}[c][0.62\textheight][c]{0.94\textwidth}
    \textbf{Placeholder for Figure 5. Subject-level confusion and error profile.}

    \vspace{0.8em}
    Panels to draw: \textbf{(a)} AIBL heldout subject-level confusion heatmap; \textbf{(b)} internal subject-level confusion heatmap; \textbf{(c)} AIBL transition/alluvial flow; \textbf{(d)} MCI-to-CN, MCI-to-AD, AD-to-MCI, AD-to-CN error subtype bars; \textbf{(e)} error-feature profile summary using MMSE, ventricle volume, hippocampus volume, AD-like z-score, confidence, and margin. The key visual claim is zero AD-to-CN errors on AIBL and boundary-like residual mistakes.
    \end{minipage}}
    \caption{\textbf{Subject-level MCI/AD error behavior.} AIBL heldout errors concentrate near adjacent disease-stage boundaries rather than collapsing disease cases into CN. Internal calibration remains weaker and should be interpreted as a limitation.}
    \label{fig:v6_error}
\end{figure}

\subsection{Algorithmic stability, risk tradeoff, and bootstrap uncertainty}
\label{sec:stability}

Bootstrap analysis with 2{,}000 resamples quantified uncertainty for the locked AIBL heldout subject-level endpoint. The 95\% CI for BAcc was 0.759--0.899, and the 95\% CI for macro AUC was 0.894--0.974. The 95\% CI for MCI recall was 0.531--0.839, and the 95\% CI for AD recall was 0.710--0.966. These intervals support the final model as the locked primary result while preserving MCI staging uncertainty as a visible limitation.

The bootstrap result also clarifies which claim is strongest. CN recall and AD-vs-CN separation are robust on AIBL heldout, while MCI recall remains the widest uncertainty interval. The paper should therefore state that MCI staging is improved and externally meaningful, but not solved.

\begin{figure}[p]
    \centering
    \fbox{\begin{minipage}[c][0.60\textheight][c]{0.94\textwidth}
    \textbf{Placeholder for Figure 6. Algorithmic ablation, calibration, risk tradeoff, and stability.}

    \vspace{0.8em}
    Panels to draw: \textbf{(a)} algorithmic ablation bars comparing best single model, arithmetic mean, equal log-pooling, partial RC-SPE, scan-level RC-SPE, and subject-level RC-SPE; \textbf{(b)} reliability curve for best single, equal log-pooling, and final RC-SPE; \textbf{(c)} risk-constraint scatterplot showing AIBL BAcc or MCI recall versus IXI false impairment rate; \textbf{(d)} leave-one-model-out BAcc stability; \textbf{(e)} bootstrap BAcc/MCI/AD recall forest plot. The visual message is that RC-SPE is an algorithmic contribution with measurable ablation, calibration, and risk-tradeoff behavior.
    \end{minipage}}
    \caption{\textbf{Algorithmic evidence and stability of RC-SPE.} Ablation, calibration, risk-tradeoff, leave-one-model-out, and bootstrap analyses show that the final locked method is not a generic average ensemble. RC-SPE improves AIBL heldout subject-level performance while preserving IXI healthy specificity and exposing MCI recall as the main residual uncertainty.}
    \label{fig:v6_stability}
\end{figure}

\subsection{Structural neurodegeneration consistency and OASIS stress test}
\label{sec:bio_results}

The AIBL heldout atlas-volume consistency score exceeded the uniform regional null in the a priori AD-key regions. The score was 0.510 versus a uniform null of 0.286, with score difference 0.225, bootstrap CI 0.479--0.526, and permutation $p=0.026$.

Across all labeled AD-relevant data, the AD-key consistency score was 0.426 versus a uniform null of 0.286, with permutation $p=0.0207$. The ADNI-only internal check remained non-significant, with score 0.342 and $p=0.1843$. These results support a bounded structural MRI neurodegeneration proxy while making clear that postmortem-stage validation is not available.

OASIS remained weak without OASIS tuning. The final subject-level model achieved OASIS accuracy 0.586, BAcc 0.334, macro AUC 0.554, AD-vs-CN AUC 0.371, and CN/MCI/AD recall 0.966/0.034/0.000. This result is reported as an unresolved transfer limitation, not as successful external validation.

\begin{figure}[p]
    \centering
    \fbox{\begin{minipage}[c][0.64\textheight][c]{0.94\textwidth}
    \textbf{Placeholder for Figure 7. CAS/Braak replacement and biological consistency.}

    \vspace{0.8em}
    Panels to draw: \textbf{(a)} old attention-only CAS failure callout; \textbf{(b)} definition of AD-key atlas regions: bilateral hippocampus, amygdala, and lateral ventricles; \textbf{(c)} AIBL heldout structural consistency score 0.510 vs null 0.286 with CI and $p=0.026$; \textbf{(d)} pooled all-labeled-AD score 0.426 vs null 0.286, plus ADNI-only non-significant internal check; \textbf{(e)} claim-boundary badges: MRI proxy only, not attention biomarker, not direct Braak proof, not clinical device.
    \end{minipage}}
    \caption{\textbf{Biological consistency and replacement of unsupported CAS/Braak claims.} The invalid attention-only CAS claim is removed. The revised biological analysis tests whether atlas structural MRI changes concentrate in a priori AD-relevant regions. The supported claim is a bounded MRI neurodegeneration proxy, not attention-map biomarker discovery or direct Braak-stage validation.}
    \label{fig:v6_biology}
\end{figure}


%% ====================================================================
\section{Discussion}
\label{sec:discussion}
%% ====================================================================

\subsection{A substantive rebuild rather than a narrow revision}
\label{sec:discussion_rebuild}

The revised study addresses the major weaknesses of the original manuscript by changing the model endpoint, evidence base, and claim boundary. Cross-dataset generalization is no longer inferred from explanation similarity. It is evaluated using a locked AIBL heldout subject-level test and an IXI healthy negative-control cohort. The unsupported attention-centered CAS claim is removed and replaced by an atlas-region structural neurodegeneration consistency analysis. The non-significant disease-stage result is no longer used to claim direct pathological staging.

This is why the revised manuscript should be read as a new evidence package rather than as a small correction. The final work asks a more defensible question: whether an atlas-guided multimodal classifier, adapted to an external cohort but evaluated on a locked heldout split, can provide strong subject-level CN/MCI/AD staging while preserving healthy specificity and exposing remaining transfer failures.

\subsection{Interpretation of the final model}
\label{sec:discussion_model}

The final model is strongest when interpreted as a domain-adapted external AD staging framework. AIBL adaptation data were used for model fitting and calibration, but AIBL heldout endpoint units remained locked and were not used for final endpoint reporting. This distinction is important. The work does not solve pure ADNI-to-AIBL zero-shot staging, but it does demonstrate that anatomically grounded MRI features and core clinical variables can support robust heldout subject-level staging within an external cohort.

The algorithmic contribution is RC-SPE rather than a new attention block. Its novelty is methodological and risk-oriented: it converts heterogeneous base-model probability streams into a locked subject-level endpoint while explicitly balancing external BAcc, MCI/AD recall, AD-to-CN error avoidance, probability calibration, and IXI healthy specificity. The ablation results show that simple averaging or equal log-pooling are insufficient, while the full subject-level RC-SPE improves AIBL BAcc and removes AD-to-CN errors. The risk-profile experiment also shows why the highest MCI-recall candidate was not selected: it increased healthy-control false impairment in IXI.

RC-SPE is therefore a pragmatic solution to a clinically important modeling problem. Single models can perform well on aggregate metrics while failing minority disease classes or producing false impairment in healthy controls. The locked algorithm explicitly balances external BAcc, MCI/AD recall, and IXI healthy specificity. Its value is therefore not only the final BAcc number, but also the combination of external heldout performance, zero IXI false impairment, bootstrap uncertainty, calibration behavior, and interpretable error behavior.

\subsection{MCI/AD boundary behavior}
\label{sec:discussion_errors}

The error analysis clarifies the clinical meaning of the remaining failures. On AIBL heldout, AD endpoint units were not missed as CN; residual AD errors were classified as MCI. MCI errors were split between correct MCI and AD, with only two MCI endpoint units classified as CN. This pattern is preferable to a model that preserves overall accuracy by collapsing minority disease classes into CN, but it also shows that precise MCI/AD boundary staging remains challenging.

MCI is biologically and clinically heterogeneous. Some MCI participants may show AD-like neurodegeneration, while others may remain stable or have non-AD etiologies. A subject-level classifier operating on retrospective structural MRI and clinical variables should therefore be expected to show uncertainty near the MCI/AD boundary. The revised manuscript makes this limitation explicit rather than hiding it behind aggregate accuracy.

\subsection{Biological consistency without overclaiming}
\label{sec:discussion_bio}

The biological analysis is intentionally bounded. The AIBL heldout atlas-volume consistency result supports disease-consistent structural MRI change in AD-relevant regions. It does not demonstrate explanation-map biomarkers, and it does not provide postmortem-stage validation. This narrower claim is more defensible and better aligned with the available data.

This change directly addresses the earlier CAS and Braak problems. The unsupported attention-only CAS is not rescued by new wording; it is removed as a central claim. The Braak claim is also not reasserted without pathology labels. Instead, the revised work provides an atlas-region MRI proxy that is statistically tested against a transparent regional null.

\subsection{OASIS and internal calibration limitations}
\label{sec:discussion_limitations}

Several limitations remain. OASIS transfer is not solved and should be treated as a stress-test failure requiring larger and cleaner external cohorts. The internal subject-level calibration pattern remains modest, especially for CN specificity within the internal split. Clinical variables contain substantial diagnostic signal, as shown by the clinical-only comparator. The 21-region atlas is anatomically interpretable but coarse, and finer parcellations may better capture cortical AD patterns.

These limitations do not invalidate the locked AIBL heldout result, but they define its correct scope. The model supports domain-adapted external heldout staging under the revised protocol. It does not support universal generalization across all public AD datasets, and it should not be described as deployment-ready clinical software.

\subsection{Open science and clinical translation boundary}
\label{sec:discussion_deployment}

The revised repository is designed for reproducible research use. It includes code for cohort construction, feature handling, model evaluation, probability ensembling, bootstrap analysis, figure generation, documentation, and a lightweight research API/frontend. This improves transparency and helps readers audit the claims. However, raw neuroimaging and clinical data remain subject to original data-use agreements and cannot be redistributed directly.

Clinically, the model should be viewed as a decision-support research prototype. Prospective multi-center validation, scanner/protocol robustness testing, local calibration, integration with clinical workflows, uncertainty reporting, and regulatory review would be required before any clinical use. The revised manuscript therefore separates deployability as software from clinical deployability as a medical product.


%% ====================================================================
\section{Conclusion}
\label{sec:conclusion}
%% ====================================================================

This revised work is a substantive rebuild of the original ARA-Net study. It replaces unsupported attention-centered claims with locked external subject-level classification, healthy negative-control specificity, bootstrap uncertainty, MCI/AD error analysis, and a bounded atlas structural neurodegeneration consistency analysis. The strongest supported claim is domain-adapted external AD staging with an MRI neurodegeneration proxy, not pure zero-shot generalization, postmortem-stage validation, OASIS success, or clinical deployment readiness.


%% ====================================================================
%% End-matter
%% ====================================================================

\section*{Data Availability}
Raw ADNI, AIBL, OASIS, and IXI data are governed by their original data-use agreements and are not redistributed. Public repository files are limited to code, aggregate reports, figures, documentation, final model configuration, and toy probability examples. Row-level subject/scan predictions, private clinical spreadsheets, MRI volumes, and model checkpoints are excluded from the public package.

\section*{Code Availability}
The public research implementation, documentation, aggregate reports, deployment wrappers, and final ensemble configuration are available at \url{https://github.com/Lava168/ARA-Net}. The released software is a research prototype; it is not a medical device and is not cleared or approved for clinical use.
